\section{Source-verified local response edits (N03)}
N02 tests a scalar decision scale; N03 instead asks where a response may change.
The anchor $b$ is our original v10 semantic response and $c$ is the fixed N01
coverage/high candidate. The same teachers and frozen geometry are reused.
Let $\delta=c-b$, $u=A\delta$, and $e=y-Ab$ on the common supported source domain.
For Huber loss $\rho$, define the integrated residual coefficient
\begin{equation}
 s_p=\begin{cases}
 [\rho(e_p)-\rho(e_p-u_p)]/u_p,&u_p\ne0,\\
 \rho'(e_p),&u_p=0.
 \end{cases}
\end{equation}
The signed and absolute allocations to element $i$ are
\begin{equation}
 g_i=\delta_i\sum_{v,p} A_{v,pi}w_{v,p}s_{v,p},\qquad
 h_i=|\delta_i|\sum_{v,p} A_{v,pi}w_{v,p}|s_{v,p}|.
\end{equation}
For nonnegative $A,w$, $|g_i|\le h_i$, and
$\sum_i g_i=\sum_{v,p}w_{v,p}[\rho(e_{v,p})-\rho(e_{v,p}-u_{v,p})]$.
This exact path accounting motivates a local gate
$q_i=\max(g_i,0)/h_i$, with zero permission for absent evidence. It is a
source-loss attribution identity, not a claim that the gated field has the same
loss reduction or higher target accuracy. A simple cancellation example shows why:
with $A=(1/2,1/2)$, $b=(0.5,0.5)$, $c=(1,0.1)$ and $y=0.55$, the candidate fits
exactly. Keeping only its positively credited first edit yields response 0.75,
which has higher source loss than the anchor's response 0.5. The quotient is an integrated loss
secant, not a newly claimed mathematical primitive.

Fit-source and held-source gates use exactly the same every-fourth eligible
source views, targets, weights and common support. The former uses the full
candidate; the latter uses a second field reconstructed without those views.
The latter gate applies to the final full-field edit only if its direction
agrees with the split-field edit; otherwise permission is zero. Geometry,
teacher tracks and the v10 anchor still use shared source information. Thus
this is a candidate-reconstruction holdout, not independent validation.

All anchored fields retain the v10 observation domain and output
$b+q\odot(c-b)$, preserving bounded values and the anchor where no update is
supported. A coverage-on-anchor-domain control isolates the domain change.
The matched uniform control separately preserves the positive and negative
Gaussian-space L1 edit masses for each query; it does not match projected
pixel counts. Six arms on all eight tasks yield 48 images sealed before target
scoring, with the unchanged 0.5 readout and both GT conventions. The held-local
arm is designated before scoring. No target RGB, GT-derived update regions,
per-task winners, or changed teachers are used.

\begin{table}[ht]
\centering\small
\resizebox{\linewidth}{!}{\begin{tabular}{lrrrrrr}
\toprule
Task & v10 & Coverage & Domain & Fit-local & Held-local & Uniform \\
\midrule
fern & 82.19 & 83.53 & 83.53 & 83.36 & 83.33 & 82.95 \\
flower & 86.53 & 87.38 & 87.38 & 87.26 & 87.25 & 86.95 \\
fortress & 97.23 & 97.46 & 97.46 & 97.42 & 97.42 & 97.53 \\
horns\_center & 95.73 & 95.55 & 95.55 & 95.89 & 95.90 & 95.85 \\
horns\_left & 92.70 & 94.43 & 94.43 & 94.16 & 94.11 & 94.00 \\
leaves & 94.08 & 95.06 & 95.06 & 95.00 & 95.00 & 95.00 \\
orchids & 68.14 & 70.05 & 70.05 & 69.64 & 69.61 & 69.04 \\
trex & 75.51 & 78.75 & 78.75 & 77.81 & 77.81 & 76.55 \\
\midrule Mean (retained GT) & 86.51 & 87.78 & 87.78 & 87.57 & 87.55 & 87.23 \\
\midrule Mean (LUDVIG GT) & 86.68 & 87.93 & 87.93 & 87.72 & 87.71 & 87.39 \\
\bottomrule
\end{tabular}
}
\caption{N03 all-eight IoU. Domain means coverage on the original anchor's
observed domain; uniform matches held-local signed edit mass.}
\label{tab:nvos-local-credit}
\end{table}
\begin{table}[ht]
\centering\small\begin{tabular}{lrrrr}
\toprule
Arm & IoU & Precision & Recall & AP \\
\midrule
v10 & 86.51 & 95.70 & 90.10 & 98.02 \\
Coverage & 87.78 & 95.83 & 91.43 & 98.15 \\
Anchor domain & 87.78 & 95.83 & 91.43 & 98.15 \\
Fit-local & 87.57 & 95.77 & 91.24 & 98.23 \\
Held-local & 87.55 & 95.76 & 91.23 & 98.27 \\
Matched uniform & 87.23 & 95.58 & 91.05 & 98.15 \\
\bottomrule
\end{tabular}

\caption{N03 macro metrics on retained GT. Every arm uses the same fixed
threshold; AP evaluates continuous scores.}
\label{tab:nvos-local-credit-macro}
\end{table}
The predeclared held-local arm obtains 87.5538\% IoU and 98.2704\% AP on retained GT.
Relative to the original v10, its mean IoU change is +1.0403 points (8 improved, 0 harmed), with a descriptive seven-scene bootstrap interval [+0.6636, +1.5292].
Relative to coverage, its mean IoU change is -0.2238 points (1 improved, 7 harmed), with a descriptive seven-scene bootstrap interval [-0.5000, -0.0479].
Relative to matched uniform attenuation, its mean IoU change is +0.3209 points (7 improved, 1 harmed), with a descriptive seven-scene bootstrap interval [+0.0750, +0.6771].
On the final fields, the common-source Huber objective is below the anchor for 8/8 held-local and 8/8 matched-uniform outputs. These fields use all source constraints, so this objective is not an independent validation score. The teacher-source criterion and target quality must be assessed separately.


The local allocation has a measured advantage over matched attenuation, but
it is not a replacement for the highest-IoU coverage candidate. Against
coverage, orchids recovers 1.4523 AP points while losing 0.4319 IoU points:
it gains 4,311 foreground pixels but loses 12,484, and adds 117 net false
positives. Trex loses 9,274 net foreground pixels and adds 9,129 net false
positives despite a lower mean background score. Relative to matched uniform
attenuation, however, the same trex local output gains 8,449 net foreground
pixels and removes 17,648 net false positives. The aggregate superiority over
uniform attenuation therefore does not establish preservation of every useful
coverage edit. All quantities are post-seal diagnostics, never update gates.

The observation-domain control changes mean IoU by only 0.000065 points.
Candidate-reconstruction holdout changes fit-local IoU by $-0.0141$ points;
we do not attribute an IoU gain to that extra reconstruction. Coverage remains
the principal best-IoU candidate, and there is no per-task selection between
coverage and local edits. Further constant or threshold searches around this
gate are not the next research step; the unresolved object extent and transport
of the admitted evidence require separate matched interventions.


\subsection{Does reference initialization discard the seed extent?}
A separate source-only audit compares the original predicted SAM reference mask
with the SAM3 track in that same authorized reference view. The fitted-query
lineage matches the declared query for every task. Agreement ranges from
96.49\% to 99.82\%; flower and orchids give 99.21\% and 97.47\% respectively.
This does not show a large reference-initialization collapse relative to the
seed. It does not establish that either mask describes the complete requested
object, or that later tracked views and geometric reconstruction are accurate.
Neither mask is ground truth; source agreement cannot substitute for target
accuracy or causally explain the target-RGB gap.
\begin{table}[ht]
\centering\small\begin{tabular}{lrrr}
\toprule
Task & Seed--track IoU & Seed retained & Track inside seed \\
\midrule
fern & 99.44 & 99.71 & 99.73 \\
flower & 99.21 & 99.57 & 99.64 \\
fortress & 99.82 & 99.92 & 99.90 \\
horns\_center & 99.56 & 99.72 & 99.83 \\
horns\_left & 99.67 & 99.83 & 99.84 \\
leaves & 99.81 & 99.88 & 99.93 \\
orchids & 97.47 & 98.37 & 99.07 \\
trex & 96.49 & 97.49 & 98.96 \\
\bottomrule
\end{tabular}

\caption{A22: agreement between two source predictions. ``Seed retained'' is
intersection/seed area, and ``track inside seed'' is intersection/track area.
These quantities are not benchmark recall or precision against ground truth.}
\end{table}
\begin{figure}[p]
\centering
\includegraphics[width=.94\linewidth,height=.87\textheight,keepaspectratio]{../../output/tmm_response_20261016/study/nvos_reference_scope_v1/reference_scope.jpg}
\caption{All eight authorized reference views, with original seed and tracked
extent. Difference panel: green overlap, blue seed-only, red track-only.
Target RGB and target GT are not used in this audit.}
\end{figure}
\IfFileExists{../../output/tmm_response_20261016/study/nvos_local_credit_spatial_v1/errors.png}{
\begin{figure}[p]
\centering
\includegraphics[width=\linewidth,height=.87\textheight,keepaspectratio]{../../output/tmm_response_20261016/study/nvos_local_credit_spatial_v1/errors.png}
\caption{N03 post-seal diagnostics on all eight tasks. Green: true positive;
red: false positive; blue: false negative. The held-local method and matched
uniform control use their predeclared rules, without task-wise selection.}
\end{figure}}{}
